\section{Introducción y repaso de la clase}

Esta es la segunda clase del curso ``Diffusion Models and Their Applications'' de KAIST CS492D, impartida por el profesor Minhyuk Sung. En esta sesión partimos desde los fundamentos de los modelos generadores introducidos en la clase anterior (Lecture 01) para ir desgranando gradualmente la idea central de los \textbf{modelos de difusión probabilística de desnaturalización} (Denoising Diffusion Probabilistic Models, DDPM).

En la clase anterior ya repasamos los siguientes conceptos clave:
\begin{itemize}
    \item Dado un conjunto de datos $\{x^{(i)}\}$, cómo considerar los datos como muestras de alguna distribución de probabilidad desconocida $p_{\text{data}}(x)$;
    \item La tarea central de los modelos generadores: aprender un mapeo desde una distribución simple (usualmente la gaussiana estándar $\mathcal{N}(0, I)$) hasta la distribución de datos;
    \item Las ideas básicas detrás de dos paradigmas de modelos generadores: GAN y VAE.
\end{itemize}

\begin{knowledgebox}{Marco unificado para los modelos generadores}
La idea central de todos los modelos generadores puede resumirse así: encontrar un mapeo desde una \textbf{distribución previa} (prior distribution) simple $p(z) = \mathcal{N}(0, I)$ hasta la \textbf{distribución de datos} (data distribution) $p_{\text{data}}(x)$. El proceso generativo consiste en: primero muestrear $z$ a partir de $p(z)$ y luego transformar $z$ en una muestra del espacio de datos $x$ mediante una función de mapeo.
\end{knowledgebox}

\subsection{Repaso de los términos clave}

Antes de entrar en DDPM, es necesario recordar los siguientes términos del marco bayesiano, los cuales atraviesan todo el curso:

\begin{table}[H]
\centering
\begin{tabular}{lll}
\toprule
\textbf{Término} & \textbf{Símbolo} & \textbf{Significado} \\
\midrule
Distribución previa (Prior) & $p(z)$ & Distribución de las variables latentes, usualmente establecida como $\mathcal{N}(0, I)$ \\
Verosimilitud (Likelihood) & $p_\theta(x|z)$ & Distribución condicional para generar datos dado una variable latente \\
Distribución marginal (Marginal) & $p_\theta(x)$ & Probabilidad marginal de los datos $\int p_\theta(x|z)p(z)dz$ \\
Distribución posterior (Posterior) & $p_\theta(z|x)$ & Distribución para inferir las variables latentes dado los datos \\
\bottomrule
\end{tabular}
\caption{Términos clave en el marco bayesiano}
\end{table}

\subsection{Resumen de la sección}

En esta sección se repasó el marco básico y los términos clave de los modelos generadores, sentando las bases para introducir posteriormente la derivación del ELBO en VAE y la expansión jerárquica de los modelos de difusión.
